\documentclass[12pt]{article}
\usepackage{amsmath,amssymb,amsfonts}
\usepackage[margin=1in]{geometry}
\begin{document}
\begin{center}
{\large\bfseries 23rd Irish Mathematical Olympiad\\24 April 2010}
\end{center}
\bigskip\noindent\textbf{Paper 1}\bigskip
\begin{enumerate}
\item Find the least $k$ for which the number $2010$ can be expressed as the sum of the squares of $k$ integers.

\item Let $ABC$ be a triangle and let $P$ denote the midpoint of the side $BC$. Suppose that there exist two points $M$ and $N$ interior to the sides $AB$ and $AC$ respectively, such that
\[
|AD| = |DM| = 2|DN|,
\]
where $D$ is the intersection point of the lines $MN$ and $AP$. Show that $|AC| = |BC|$.

\item Suppose $x$, $y$, $z$ are positive numbers such that $x + y + z = 1$. Prove that

\begin{enumerate}
\item $xy + yz + zx \ge 9xyz$;

\item $xy + yz + zx < \dfrac{1}{4} + 3xyz$.
\end{enumerate}

\item The country of Harpland has three types of coin: \textit{green}, \textit{white} and \textit{orange}. The unit of currency in Harpland is the shilling. Any coin is worth a positive integer number of shillings, but coins of the same colour may be worth different amounts. A set of coins is stacked in the form of an equilateral triangle of side $n$ coins, as shown below for the case of $n = 6$.

The stacking has the following properties:
\begin{enumerate}
\item no coin touches another coin of the same colour;
\item the total worth, in shillings, of the coins lying on any line parallel to one of the sides of the triangle is divisible by three.
\end{enumerate}
Prove that the total worth in shillings of the \textit{green} coins in the triangle is divisible by three.

\item Find all polynomials $f(x) = x^3 + bx^2 + cx + d$, where $b$, $c$, $d$ are real numbers, such that $f(x^2 - 2) = -f(-x)f(x)$.
\end{enumerate}

\newpage
\begin{center}
{\large\bfseries 23rd Irish Mathematical Olympiad\\24 April 2010}
\end{center}
\bigskip\noindent\textbf{Paper 2}\bigskip
\begin{enumerate}
\setcounter{enumi}{5}
\item There are $14$ boys in a class. Each boy is asked how many other boys in the class have his first name, and how many have his last name. It turns out that each number from $0$ to $6$ occurs among the answers.

Prove that there are two boys in the class with the same first name and the same last name.

\item For each odd integer $p \ge 3$ find the number of real roots of the polynomial
\[
f_p(x) = (x-1)(x-2)\cdots(x-p+1) + 1.
\]

\item In the triangle $ABC$ we have $|AB| = 1$ and $\angle ABC = 120^\circ$. The perpendicular line to $AB$ at $B$ meets $AC$ at $D$ such that $|DC| = 1$. Find the length of $AD$.

\item Let $n \ge 3$ be an integer and $a_1, a_2, \ldots, a_n$ be a finite sequence of positive integers, such that, for $k = 2, 3, \ldots, n$
\[
n(a_k + 1) - (n-1)a_{k-1} = 1.
\]
Prove that $a_n$ is not divisible by $(n-1)^2$.

\item Suppose $a$, $b$, $c$ are the side lengths of a triangle $ABC$. Show that
\[
x = \sqrt{a(b+c-a)}, \quad y = \sqrt{b(c+a-b)}, \quad z = \sqrt{c(a+b-c)}
\]
are the side lengths of an acute-angled triangle $XYZ$, with the same area as $ABC$, but with a smaller perimeter, unless $ABC$ is equilateral.
\end{enumerate}
\end{document}
